\chapter{Ethics and Data Management}
\label{edm}

% =============================================================================
% ETHICS
% =============================================================================

\section{Ethical Considerations}

% TODO: Acknowledge UvA ethical code and RDM policies
%   "I acknowledge that the thesis adheres to the ethical code and
%    research data management policies of UvA and IvI."

% TODO: Ethical considerations specific to this project:
%   - LLM agents are NOT human subjects — no IRB/ethics board needed
%   - But: LLM behavior reflects training data biases (cooperation/conflict
%     patterns may encode cultural biases from pretraining corpus)
%   - Social simulation implications: if LLM-ABMs are used for policy,
%     reasoning depth effects must be accounted for (our finding that
%     deeper reasoning amplifies structural incentives has implications)
%   - No personal data involved — all agents are synthetic
%   - Compute footprint: ~1000+ SBU on Snellius H100 GPUs. Acknowledge
%     environmental cost of large-scale LLM inference experiments.

% =============================================================================
% AI TOOL USAGE STATEMENT
% =============================================================================

\section{AI Tool Usage}
\label{sec:ai-usage}

% TODO: Write transparent disclosure of Claude Code usage during this project.
%   This section should be honest and specific. Koen writes this himself.
%
% Structure to follow:
%
% 1. WHAT TOOL
%   - Claude Code (Anthropic's CLI agent, Claude Opus 4.6 model)
%   - Used throughout the project (Feb-Jul 2026)
%   - Governed by a self-imposed "Thesis Assistance Manifesto" (CLAUDE.md)
%     that defined boundaries for acceptable vs. unacceptable assistance
%
% 2. WHAT IT WAS USED FOR (be specific)
%   - Code development: game engine, sweep infrastructure, analysis pipelines,
%     prompt engineering for reasoning depth manipulation
%   - Experiment orchestration: Snellius job submission scripts, GPU scheduling,
%     auto-submit scripts for parameter sweeps
%   - Data analysis: metrics computation, statistical analysis scripts,
%     trace analysis pipeline
%   - Literature research: finding relevant papers, explaining technical concepts,
%     connecting ideas across sources
%   - Project management: sprint planning, roadmap updates, meeting preparation
%   - Thesis structure: chapter outlines (TODO comments only, not prose)
%   - Debugging: code review, error diagnosis, infrastructure troubleshooting
%
% 3. WHAT IT WAS NOT USED FOR (equally important)
%   - NO thesis prose was generated by Claude — all text written by Koen
%   - NO research decisions were made by Claude — all design choices are Koen's
%   - NO results interpretation was delegated — all analysis conclusions are Koen's
%   - NO literature was cited without Koen reading the primary source
%   - Claude was explicitly prohibited from writing prose, making decisions,
%     or replacing Koen's thinking (enforced via manifesto)
%
% 4. HOW QUALITY WAS MAINTAINED
%   - Self-imposed manifesto with strict rules (see Appendix \ref{AppendixManifesto})
%   - "Explanation Test": Koen must be able to explain every concept, defend
%     every choice, and modify any code without assistance
%   - "Voice Test": thesis must sound like Koen, not like an LLM
%   - Regular work sessions without Claude to verify independent capability
%   - All code understood line-by-line before inclusion
%   - Supervisor (Debraj Roy) informed of AI usage throughout
%
% 5. REFLECTION
%   - What worked well about this approach?
%   - What were the risks / moments of over-reliance?
%   - What would you do differently?
%   - How does this relate to the broader question of AI in academic work?
%
% 6. CITE UvA AI POLICY
%   - Reference UvA's policy on AI tools in academic work
%   - Confirm compliance with institutional requirements
%
% NOTE TO KOEN: This is the ONE section where you should be most honest and
% reflective. Reviewers and examiners will respect transparency. The manifesto
% (Appendix) shows you thought carefully about this. Don't undersell the
% assistance, don't oversell it. Just be precise about what happened.

% =============================================================================
% DATA MANAGEMENT
% =============================================================================

\section{Data Management}

% TODO: Fill in actual data table with:

\begin{table}[h]
\centering
\begin{tabular}{|l|l|l|}
\hline
\textbf{\begin{tabular}[c]{@{}l@{}}Short description\\(max. 10 words)\end{tabular}} &
\textbf{\begin{tabular}[c]{@{}l@{}}Availability\\(e.g., URL, DOI)\end{tabular}} &
\textbf{\begin{tabular}[c]{@{}l@{}}License\\(e.g., MIT, GPL, CC)\end{tabular}} \\ \hline

% TODO: Simulation source code
% Origins simulation engine & GitHub URL & MIT \\ \hline

% TODO: Experiment configurations
% Experiment YAML configs & GitHub URL & MIT \\ \hline

% TODO: Run data (history, traces, metrics)
% Simulation run data (~X GB) & Zenodo DOI / Figshare & CC-BY-4.0 \\ \hline

% TODO: Analysis scripts
% Analysis and plotting scripts & GitHub URL & MIT \\ \hline

% TODO: Model weights (not ours, but document)
% Gemma 2 27B IT model weights & HuggingFace & Gemma License \\ \hline

% TODO: Thesis LaTeX source
% Thesis LaTeX source & GitHub URL & CC-BY-4.0 \\ \hline

\end{tabular}
\caption{Data and code artifacts for this thesis.}
\label{tab:data-management}
\end{table}

% TODO: Reproducibility statement
%   - All runs have unique seeds, logged configs, model version, hardware specs
%   - Full prompt text versioned in repository
%   - Scripts to reproduce every figure from raw data
%   - See publishable_checklist.md Section 8 for full requirements
